% ============================================================================
% epsilon_intro.tex — Competitive draft (regulatory-economics framing)
% "The Economics of Extended Play: What Happens When a Regulator
%  Adds Time to the Clock?"
% ============================================================================

% ----------------------------------------------------------------------------
%  ABSTRACT
% ----------------------------------------------------------------------------
\begin{abstract}

How do enforcers, participants, and markets respond to a regulatory mandate
that alters the basic parameters of a competitive environment?
We study a 2024 FIFA directive requiring referees to add more stoppage time
to professional football matches---a setting that provides unusually clean
identification: a discrete regulatory shock, individually observable
enforcer compliance, high-frequency behavioral data on regulated agents,
and real-time market prices.
Using difference-in-differences across 36{,}267 matches in 15 country-leagues,
we estimate that the directive increased match duration by
$\hat{\beta}=0.744$ minutes (wild cluster bootstrap $p=0.036$).
Five findings emerge.
First, enforcers complied near-universally: every referee with at least
20 post-treatment matches shifted behavior, a 100\% compliance rate
unusual in regulatory settings.
Second, regulated agents adapted immediately, reducing their per-minute
goal-scoring intensity by 31\%.
Third, this behavioral offset was exact---total goals per match were
unchanged ($\hat{\beta}=-0.033$, $p=0.39$).
Fourth, betting markets showed no forecast deterioration
(Pinnacle Brier-score DiD: $+0.007$, $p=0.41$), indicating that
information processors absorbed the structural change without friction.
Fifth, we find no adverse welfare effects: injuries declined
($-0.090$, $p=0.005$) and muscular injuries were unchanged ($p=0.73$).
Pre-treatment trends are not parallel (F$=9.0$, $p<0.001$), and the
\citet{RambachanRoth2023} honest-DiD sensitivity bound is
$\bar{M}^{*}=0$; we address this with permutation inference ($p=0.045$)
and a 420-specification curve (median $\hat{\beta}=0.669$, 88.8\%
of specifications significant).

\end{abstract}

% ----------------------------------------------------------------------------
%  INTRODUCTION
% ----------------------------------------------------------------------------
\section{Introduction}\label{sec:intro}

When a regulator changes the rules, how completely do enforcers comply?
How quickly do regulated agents adapt?
And do the markets that price the regulated activity recalibrate, or do
they stumble?
These questions sit at the core of regulatory economics, yet clean
empirical answers are scarce.
Compliance is hard to observe at the individual-enforcer level.
Behavioral adaptation unfolds over months or years, confounded by
concurrent shocks.
Market responses are entangled with information asymmetries.
We study a setting in which all three channels---compliance, adaptation,
and market absorption---are observable at high frequency, with a sharp
regulatory discontinuity and thousands of post-treatment observations.

The setting is professional football (soccer).
In June 2024, FIFA issued a directive instructing referees worldwide to
enforce existing stoppage-time rules more strictly, effectively adding
minutes to every match.
The directive was not a law change.
The rules governing stoppage time had been on the books for decades.
What changed was regulatory emphasis: FIFA signaled, through public
communications and referee briefings, that under-enforcement would no
longer be tolerated.
This creates a natural experiment.
Treated leagues---those under strong FIFA compliance pressure---adopted
the directive; control leagues did not, or adopted it later.
The treatment is discrete, its timing is known, and the outcome
(match duration) is measured to the second.
Our preferred difference-in-differences specification yields
$\hat{\beta}=0.744$ additional minutes of stoppage time
(wild cluster bootstrap $p=0.036$, 15 clusters, $N=36{,}267$).
Within-country estimates are similar ($\hat{\beta}=0.868$, WCB $p=0.074$,
8 clusters).
Restricting to the Big~5 European leagues sharpens the effect:
$\hat{\beta}=1.575$ ($p<0.001$).

The compliance result is striking.
Of 85 referees with at least 20 post-directive matches, every one
shifted behavior in the direction the regulator intended---a 100\%
compliance rate.
This is unusual.
The regulatory-compliance literature documents wide variation in enforcer
behavior, even under direct monitoring
\citep{DufloGreenstoneHanna2013,FinanOlken2017}.
Referees are quasi-autonomous agents operating in real time, under
intense public scrutiny, with no formal penalty for non-compliance.
Yet compliance was both universal and immediate: a
\citet{Gelbach2016} decomposition attributes 87\% of the
post-treatment increase in late goals to quantity (more time added)
rather than intensity (more goals per minute of added time),
and 87\% of the treatment effect materialized within the first
five matchweeks.
The mechanism appears to be coordination on a new norm rather than
coercion---a finding that speaks to the literature on soft regulation
and focal-point shifts \citep{Sunstein1996,McAdams2000}.

The behavioral-adaptation result is equally clean.
If referees add time and players maintain their pre-treatment intensity,
total goals should rise mechanically.
They do not.
The per-minute goal rate fell from 0.72 to 0.50---a 31\% decline---and
total goals per match were unchanged
($\hat{\beta}=-0.033$, $p=0.39$).
The offset is almost exactly one-for-one.
This is a textbook \citet{Peltzman1975} result: a safety or structural
regulation intended to shift one margin provokes a compensating response
on another.
Peltzman's original finding---that mandatory seatbelts reduced driver
deaths but increased pedestrian deaths---has been debated for decades.
Our setting offers a cleaner test because the ``risk compensation'' is
directly measurable.
Players did not merely adjust vaguely; they reduced scoring intensity by
precisely the amount needed to leave the aggregate outcome unchanged.
Goals were displaced, not created: the coefficient on goals scored after
the 90th minute is $+0.030$ ($p<0.001$), and the Gelbach decomposition
confirms that 87\% of this displacement reflects additional time
rather than heightened late-game intensity.

We make five contributions.
%
\begin{enumerate}[label=(\roman*)]

\item \textbf{Enforcer compliance.}
We document universal compliance with a voluntary regulatory directive,
at the individual-enforcer level, in a sample of 85 referees.
This contributes to the literature on regulatory enforcement and
inspections \citep{DufloGreenstoneHanna2013,DimantTosato2018}.

\item \textbf{Behavioral adaptation.}
We show that regulated agents offset the treatment effect completely,
reducing per-minute intensity by 31\% so that total output is unchanged.
This is among the cleanest demonstrations of the Peltzman offset
in a non-laboratory setting \citep{Peltzman1975,VisscherSi1997}.

\item \textbf{Market absorption.}
Betting markets---real-time information processors with strong
incentives for accuracy---show no forecast deterioration.
Pinnacle's Brier-score DiD is $+0.007$ ($p=0.41$); Bet365's is
$+0.005$ ($p=0.58$).
This contributes to the efficient-markets literature in prediction
settings \citep{Thaler1988,GilMollaPazZone2019}.

\item \textbf{Welfare neutrality.}
Total injuries declined ($-0.090$, $p=0.005$).
Muscular injuries---the most plausible adverse channel---were
unchanged ($p=0.73$), and the estimate is well-powered.
The directive achieved its compliance objective without detectable
welfare costs.

\item \textbf{Identification under imperfect pre-trends.}
Our pre-treatment trends are not parallel ($F=9.0$, $p<0.001$),
and the \citet{RambachanRoth2023} sensitivity analysis yields
$\bar{M}^{*}=0$.
We address this transparently rather than ignoring it, and show
that the result survives permutation inference ($p=0.045$), a
420-specification curve (median $\hat{\beta}=0.669$,
88.8\% significant), and leave-one-out analysis
(range $[0.587, 0.972]$, all significant).
Excluding the 2022--23 season strengthens the estimate
($\hat{\beta}=0.810$).

\end{enumerate}

We are candid about what identification can and cannot deliver here.
The pre-trend failure means that a pure parallel-trends defense is
unavailable.
The \citet{RambachanRoth2023} HonestDiD analysis confirms this:
the identified set includes zero at $\bar{M}^{*}=0$, the weakest
possible sensitivity bound.
Two considerations temper this concern.
First, the permutation test ($p=0.045$) is agnostic to parallel
trends; it asks only whether the treated-control contrast is unusually
large relative to placebo assignments.
Second, the specification curve spans 420 defensible specifications---varying
the control group, functional form, and covariate set---and 88.8\%
reject the null.
The leave-one-out range is tight ($[0.587, 0.972]$), no single
league drives the result, and excluding pre-treatment seasons that
contribute most to the trend break strengthens the estimate.
We present these results as a package. No single robustness check
is dispositive; together, they support a treatment effect in the
range of 0.6--1.0 additional minutes.

This paper contributes to several literatures.
Most directly, it speaks to the economics of regulatory compliance.
\citet{DufloGreenstoneHanna2013} show that third-party auditors
in India's pollution-monitoring regime can be induced toward honest
reporting through structural reform; we document a case where
enforcers shifted behavior through soft coordination alone.
Second, the Peltzman offset we identify connects to a long line of
work on behavioral responses to regulation
\citep{Peltzman1975,BlomquistPeltzman1986,CohenEinaav2003},
with the distinction that our offset is exact and directly measured.
Third, the betting-market results contribute to the literature on
prediction-market efficiency
\citep{Thaler1988,SnowbergWolfers2010,BordaloGennaioli2019},
showing that a structural rule change does not degrade forecast
quality.
Finally, the paper sits within the growing literature that uses
sports as a laboratory for economic questions
\citep{ChiapporiGrosecloseLevitt2002,RomerBR2006,GarrechtGaricano2006},
with the advantage that our regulatory shock is discrete and our
outcomes are measured with unusual precision.

The remainder of the paper proceeds as follows.
Section~\ref{sec:background} describes the FIFA directive and the
institutional setting.
Section~\ref{sec:data} introduces the data.
Section~\ref{sec:empirical} presents the empirical strategy.
Section~\ref{sec:results} reports the main results.
Section~\ref{sec:mechanisms} investigates compliance, behavioral
adaptation, and market absorption.
Section~\ref{sec:robustness} addresses robustness.
Section~\ref{sec:conclusion} concludes.
